\section{Hall-Effect}
\begin{frame} \ftx{Hall-Effekt}
    \s{
        The Hall effect describes the occurrence of an electrical voltage within a current-carrying conductor due to the displacement of charge by a magnetic field. \par
        In Figure \ref{Hall-Effekt}, the conductor is traversed by a magnetic field, similar to the experimental setup for the Lorentz force. However, the conductor is shown here in plate form so that the charge carriers can move across the entire surface of the conductor.        
        When current flows through the conductor, the Lorentz force acts on the moving charges. According to the right-hand rule, this causes a perpendicular deflection of the charge carriers in the conductor towards the edge. The charge displacement creates an electric field in the conductor that is opposite to the Lorentz force. The voltage difference resulting from the charge displacement can be measured at the edges of the conductor. It is referred to as the Hall voltage $U_{\mathrm{H}}$.
    }

    \begin{columns}
        \renewcommand{\figurename}{Figure}
        \column[c]{0.5\textwidth}
        
            \b{
                \only<1>{
                    \foo{}{width=\textwidth}{Hall-EffektA}{
                    \put(7, 53.5){\color{current}$I$}
                    \put(81, 11){\color{current}$I$}
                    }{}
                }
                \only<2>{
                    \foo{}{width=\textwidth}{Hall-EffektB}{
                    \put(7, 53.5){\color{current}$I$}
                    \put(81, 11){\color{current}$I$}
                    \put(37, 69){\color{magnetfeld}$B$}
                    }{}
                }
                \only<3>{
                \foo{}{width=\textwidth}{Hall-EffektC}{
                    \put(7, 53.5){\color{current}$I$}
                    \put(81, 11){\color{current}$I$}
                    \put(91, 30.5){\color{gray}$d$}
                    \put(37, 69){\color{magnetfeld}$B$}
                    \put(53, 52){\color{voltage}$U_H$}
                    }{}
                }
            }
            
            \only<4->{
                \fo{width=0.55\textwidth}{width=\textwidth}{Hall-EffektD}{
                    \put(7, 53.5){\color{current}$I$}
                    \put(81, 11){\color{current}$I$}
                    \put(91, 30.5){\color{gray}$d$}
                    \put(37, 69){\color{magnetfeld}$\vec{B}$}
                    \put(53, 52){\color{voltage}$U_H$}
                }{{\bf Experimental setup for the Hall effect.} The Lorentz force generated by the magnetic field deflects the moving charge carriers perpendicularly to the edge. This charge displacement generates an electrically measurable voltage - the Hall voltage. \label{Hall-Effekt}}
            }

        \column[c]{0.5\textwidth}
            \b{
                \begin{itemize}
                    \item<1-> A current  $I$ flows through a conductor plate.
                    \item<2-> The Lorentz force deflects the moving charge carriers to the edge.
                    \item<3-> The charge displacement generates an electric field.
                    \item<4-> This is measured as the Hall voltage $U_{\mathrm{H}}$.
                \end{itemize}
            
                \onslide<4->{
                    \begin{eq}
                        U_{\mathrm{H}}  = A_{\mathrm{H}} \cdot \frac{I \cdot B}{d}
                    \end{eq}
                }
            }
    \end{columns}
    \s{
        The Hall constant $A_\mathrm{H}$ is required to determine the Hall voltage $U_\mathrm{H}$. The Hall constant $A_\mathrm{H}$ is a material parameter that provides information about the charge carrier density and thus about the conductivity of the material.
         
        The Hall voltage can be calculated (as shown in equation \ref{GlHallSpannung}) by multiplying the Hall constant $A_{\mathrm{H}}$ by the product of the current $I$ and the magnetic flux density $B$ divided by the thickness of the current-carrying plate $d$.

        \begin{equation}
            U_{\mathrm{H}} = A_{\mathrm{H}} \cdot \frac{I \cdot B}{d} \label{GlHallSpannung}
        \end{equation}\\
        \begin{Key point}{Hall-Effect}
            The Hall voltage $U_{\mathrm{H}}$ describes the occurrence of an electrical voltage within the current-carrying conductor perpendicular to the current flow due to an acting magnetic field.
        \end{Key point}
    }
\end{frame}


\begin{frame} \ftx{Beispiel: Hall-Spannung}
    \begin{expl}{Hall-Voltage}{}
        A metal plate with a thickness of $d = 1\,\mathrm{cm}$ and a current of $I = 2\,\mathrm{A}$ is located in a magnetic field with a flux density of $B = 0.5\,\mathrm{T}$. The Hall constant of the material is $A_{\mathrm{H}} = 3.2 \cdot 10^{-3} \, \frac{\mathrm{m}^3}{\mathrm{C}}$. 
        Calculate the Hall voltage.
        
        \begin{align*}
            \onslide<2->{U_{\mathrm{H}} & = A_{\mathrm{H}} \cdot \frac{I \cdot B}{d}}\\
            \onslide<3->{U_{\mathrm{H}} & = 3,2 \cdot 10^{-3} \, \tfrac{\mathrm{m}^3}{\mathrm{C}} \cdot \frac{2 \, \mathrm{A} \cdot 0,5 \, \mathrm{T}}{0,01 \, \mathrm{m}}}\\
            \onslide<4->{U_{\mathrm{H}} & = 0,32 \, \mathrm{V}}
        \end{align*}
        
    \end{expl}
\end{frame}
